\documentclass[11pt]{article}
\usepackage[margin=1in]{geometry}
\usepackage{amsmath,amssymb,amsthm}
\usepackage{booktabs}
\usepackage{graphicx}
\usepackage{microtype}
\usepackage{xcolor}
\usepackage[colorlinks=true,linkcolor=blue!45!black,citecolor=blue!45!black,urlcolor=blue!45!black]{hyperref}

\newtheorem{theorem}{Theorem}
\newtheorem{lemma}{Lemma}
\newcommand{\Sn}{S_n}
\newcommand{\Fix}{\operatorname{Fix}}
\newcommand{\ord}{\operatorname{ord}}

\title{Exact Enumeration and Scalar Recurrences for\\Left-Right Orbits of Subsets of Permutations}
\author{Candidate full solution for arXiv:2503.09542}
\date{26 August 2026}

\begin{document}
\maketitle

\begin{center}
\fcolorbox{blue!45!black}{blue!4}{\parbox{0.91\textwidth}{
\textbf{Status: candidate full solution, likely valid.}
The formula below determines every $f_{n,k}$ in Question 3.5 of
Kushwaha--Tripathi~\cite{KT}.  A constructive Wronskian elimination then
gives a scalar finite-shift recurrence in $k$ for every fixed $n$.
}}
\end{center}

\section{Source question}

Kushwaha and Tripathi define $\mathcal P_{n,k}=\{A\subseteq S_n:|A|=k\}$
and declare $A\sim B$ if $B=\mu A\nu$ for some $\mu,\nu\in S_n$.
Their Question 3.5 (PDF page 12, source label \texttt{ques:Equivalence}) is:

\begin{quote}
Determine the number of equivalence classes in $\mathcal P_{n,k}$ under
$\sim$, that is, determine
$f_{n,k}=|\mathcal P_{n,k}/\!\sim|$. Does $f_{n,k}$ satisfy some recurrence?
\end{quote}

\begin{figure}[ht]
  \centering
  \includegraphics[width=\textwidth]{figures/open_problem_crop.png}
  \caption{The definition and Question 3.5 in the source paper, PDF page 12.}
\end{figure}

\section{Notation and explicit formula}

Write a partition $\lambda\vdash n$ as
$\lambda=(1^{m_1}2^{m_2}\cdots)$ and put
\[
 z_\lambda=\prod_{j\geq 1}j^{m_j}m_j!.
\]
Thus $z_\lambda$ is the centralizer size of a permutation of type $\lambda$.
For $r\geq1$, let $\lambda^{[r]}$ be the cycle type of the $r$-th power
of such a permutation. Explicitly, every part $j$ of $\lambda$ is replaced
by $\gcd(j,r)$ parts of size $j/\gcd(j,r)$. Also set
\[
 o(\lambda)=\mathop{\mathrm{lcm}}\{j:m_j>0\}.
\]

For $\lambda,\eta\vdash n$, let
$L_{\lambda,\eta}=\mathrm{lcm}(o(\lambda),o(\eta))$. For every
$r\mid L_{\lambda,\eta}$ define
\[
 a_r(\lambda,\eta)=
 \begin{cases}
 z_{\lambda^{[r]}},&\lambda^{[r]}=\eta^{[r]},\\
 0,&\lambda^{[r]}\neq\eta^{[r]},
 \end{cases}
\]
and, for every $\ell\mid L_{\lambda,\eta}$, define
\[
 c_\ell(\lambda,\eta)
 =\frac1\ell\sum_{r\mid\ell}
 \mu_{\mathrm M}(\ell/r)a_r(\lambda,\eta),
\]
where $\mu_{\mathrm M}$ is the number-theoretic Moebius function.

\begin{theorem}[Exact orbit formula]
For every $n\geq1$ and $0\leq k\leq n!$,
\begin{equation}\label{eq:main}
\boxed{
f_{n,k}=
\sum_{\lambda,\eta\vdash n}\frac1{z_\lambda z_\eta}
[x^k]\prod_{\ell\mid L_{\lambda,\eta}}
(1+x^\ell)^{c_\ell(\lambda,\eta)}.}
\end{equation}
In particular, the right-hand side is a finite explicit determination of
$f_{n,k}$ in terms of partitions of $n$.
\end{theorem}

\section{Proof intuition}

Replace the paper's right multiplier by its inverse, which does not change
the equivalence relation.  A pair $(\mu,\nu)$ then permutes $S_n$ by
$\sigma\mapsto\mu\sigma\nu^{-1}$. A subset is fixed precisely when it is a
union of cycles of this permutation, so its size-generating polynomial is a
product of factors $1+x^{\text{cycle length}}$.

The only non-obvious input is that all of those cycle lengths can be recovered
without iterating through $S_n$. The fixed points of the $r$-th power are
exactly the conjugators between $\nu^r$ and $\mu^r$. Their number is either
zero or the centralizer size of their common power type. Moebius inversion
turns these fixed-point counts into exact action-cycle counts. Burnside's
lemma then averages the fixed-subset polynomial, and grouping the average by
the two conjugacy types gives~\eqref{eq:main}.

For the recurrence clause, regard the finitely many fixed-subset polynomials
as an explicitly known polynomial space. A basis and its Wronskian produce a
scalar differential equation annihilating every polynomial in that space,
including $F_n(x)=\sum_k f_{n,k}x^k$. Coefficient extraction gives a recurrence
for $f_{n,k}$ whose coefficients are computed only from the partition data,
not from a previously expanded row of orbit counts.

\section{Proof}

For $\mu,\nu\in S_n$, define a permutation of the set $S_n$ by
\[
 T_{\mu,\nu}(\sigma)=\mu\sigma\nu^{-1}.
\]
These transformations give the usual $S_n\times S_n$ action, and their
orbits on $k$-subsets are exactly the equivalence classes in the question.

Suppose $\mu$ and $\nu$ have cycle types $\lambda$ and $\eta$. Since
\[
 T_{\mu,\nu}^{r}(\sigma)=\mu^r\sigma\nu^{-r},
\]
we have
\begin{align*}
 \sigma\in\Fix(T_{\mu,\nu}^{r})
 &\Longleftrightarrow \mu^r\sigma\nu^{-r}=\sigma\\
 &\Longleftrightarrow \mu^r=\sigma\nu^r\sigma^{-1}.
\end{align*}
Thus there are no fixed points unless
$\lambda^{[r]}=\eta^{[r]}$. If the two power types agree, the set of
conjugators is a coset of the centralizer of $\nu^r$, and hence
\begin{equation}\label{eq:fixpower}
 |\Fix(T_{\mu,\nu}^{r})|=a_r(\lambda,\eta).
\end{equation}

Let $C_\ell$ denote the number of $\ell$-cycles of $T_{\mu,\nu}$. Its
order divides $L_{\lambda,\eta}$. For each $r\mid L_{\lambda,\eta}$,
a cycle of length $\ell$ contributes all of its $\ell$ points to
$\Fix(T^r)$ exactly when $\ell\mid r$. Therefore
\[
 a_r(\lambda,\eta)=\sum_{\ell\mid r}\ell C_\ell.
\]
Moebius inversion gives
\[
 C_\ell=\frac1\ell\sum_{r\mid\ell}
 \mu_{\mathrm M}(\ell/r)a_r(\lambda,\eta)
 =c_\ell(\lambda,\eta).
\]
This also proves directly that every $c_\ell(\lambda,\eta)$ defined above
is a nonnegative integer.

A subset of $S_n$ is fixed by $T_{\mu,\nu}$ if and only if it is a union
of cycles of $T_{\mu,\nu}$. Consequently, the number of its fixed subsets
of size $k$ is
\begin{equation}\label{eq:fixedsubset}
 [x^k]\prod_{\ell\mid L_{\lambda,\eta}}
 (1+x^\ell)^{c_\ell(\lambda,\eta)}.
\end{equation}

Burnside's lemma~\cite{Cameron} averages~\eqref{eq:fixedsubset} over the
$(n!)^2$ pairs $(\mu,\nu)$. There are $n!/z_\lambda$ permutations of type
$\lambda$ and $n!/z_\eta$ permutations of type $\eta$. Grouping the
average by $(\lambda,\eta)$ therefore replaces the multiplicity divided by
$(n!)^2$ by $1/(z_\lambda z_\eta)$, which yields~\eqref{eq:main}. \qed

\section{The scalar recurrence}

Equation~\eqref{eq:main} gives the whole row generating polynomial
\[
 F_n(x)=\sum_{k=0}^{n!}f_{n,k}x^k
 =\sum_{\lambda,\eta\vdash n}\frac1{z_\lambda z_\eta}
 \prod_{\ell\mid L_{\lambda,\eta}}
 (1+x^\ell)^{c_\ell(\lambda,\eta)}.
\]
It immediately provides an exact algorithm indexed by $p(n)^2$ pairs of
partitions, rather than by the $\binom{n!}{k}$ subsets in the original
definition. It also recovers the source's complement symmetry
$f_{n,k}=f_{n,n!-k}$, because every factor in a fixed-subset polynomial is
palindromic of total degree $n!$.

Set
\[
 Q_{\lambda,\eta}(x)=
 \prod_{\ell\mid L_{\lambda,\eta}}
 (1+x^\ell)^{c_\ell(\lambda,\eta)}
\]
and let $V_n$ be the rational linear span of all these polynomials. Since
every $Q_{\lambda,\eta}$ is given explicitly by cycle-type data, a basis
$Q_1,\ldots,Q_{m_n}$ of $V_n$ is obtained constructively by row reduction
of their coefficient vectors.

For $0\leq j\leq m_n$, define
\begin{equation}\label{eq:Aj}
A_{n,j}(x)=(-1)^{m_n+j}
\det\!\left(Q_i^{(r)}(x)\right)_{
\substack{0\leq r\leq m_n,\ r\neq j\\1\leq i\leq m_n}}.
\end{equation}
In particular, $A_{n,m_n}$ is the Wronskian of the chosen basis and is
nonzero over the characteristic-zero field $\mathbb Q$~\cite{Ince}.

\begin{theorem}[Constructive scalar recurrence]\label{thm:recurrence}
For every fixed $n$, the row generating polynomial satisfies
\begin{equation}\label{eq:ode}
 \sum_{j=0}^{m_n}A_{n,j}(x)F_n^{(j)}(x)=0.
\end{equation}
Write $A_{n,j}(x)=\sum_t a_{n,j,t}x^t$, extend $f_{n,q}=0$ outside
$0\leq q\leq n!$, and write
$q^{\underline{j}}=q(q-1)\cdots(q-j+1)$. Then, for every $r\geq0$,
\begin{equation}\label{eq:scalarrecurrence}
\boxed{
\sum_{j=0}^{m_n}\sum_t
a_{n,j,t}(r+j-t)^{\underline{j}}f_{n,r+j-t}=0.}
\end{equation}
The polynomials $A_{n,j}$, hence all coefficients in this scalar recurrence,
are computed from the partition power-types in the orbit formula without
using the values of $f_{n,k}$.
\end{theorem}

\begin{proof}
Append a last column consisting of $y,y',\ldots,y^{(m_n)}$ to the derivative
matrix of $Q_1,\ldots,Q_{m_n}$. Expanding its determinant along that last
column gives the left side of~\eqref{eq:ode}. If $y\in V_n$, the columns are
linearly dependent, so the determinant vanishes. Formula~\eqref{eq:main}
places $F_n$ in $V_n$, proving~\eqref{eq:ode}. The coefficient of $x^r$ in
$x^tF_n^{(j)}(x)$ is
$(r+j-t)^{\underline{j}}f_{n,r+j-t}$. Extracting $[x^r]$ proves
\eqref{eq:scalarrecurrence}.
\end{proof}

There is also a lower-dimensional practical recurrence before scalar
elimination. If
$q_{\lambda,\eta,k}=[x^k]Q_{\lambda,\eta}(x)$, put
\[
b_{\lambda,\eta,s}=
\sum_{\ell\mid s}(-1)^{s/\ell-1}\ell c_\ell(\lambda,\eta).
\]
The logarithmic derivative of $Q_{\lambda,\eta}$ gives
\begin{equation}\label{eq:auxrecurrence}
q_{\lambda,\eta,0}=1,\qquad
kq_{\lambda,\eta,k}=
\sum_{s=1}^{k}b_{\lambda,\eta,s}q_{\lambda,\eta,k-s},
\end{equation}
and the desired scalar count is the weighted sum
$f_{n,k}=\sum_{\lambda,\eta\vdash n}(z_\lambda z_\eta)^{-1}
q_{\lambda,\eta,k}$. Thus both the scalar existence statement and an
efficient recursive computation are explicit.

\subsection*{Example: the complete recurrence for $n=3$}

For $n=3$, the six distinct fixed-subset polynomials span a
four-dimensional space. Wronskian elimination gives
\begin{align*}
0={}&(x^2-1)^3F_3^{(4)}-12x(x^2-1)^2F_3^{(3)}
+12(x^2-1)(6x^2-1)F_3''\\
&-120x(2x^2-1)F_3'+360x^2F_3.
\end{align*}
Equivalently, with $f_{3,q}=0$ outside $0\leq q\leq6$, for every $r\geq0$,
\begin{align*}
0={}&-(r+1)(r+2)(r+3)(r+4)f_{3,r+4}\\
&+3(r-4)(r-1)(r+1)(r+2)f_{3,r+2}\\
&-3r(r-6)(r-5)(r-3)f_{3,r}\\
&+(r-8)(r-7)(r-6)(r-5)f_{3,r-2}.
\end{align*}
Together with $(f_{3,0},f_{3,1},f_{3,2},f_{3,3})=(1,1,2,2)$, the cases
$r=0,1,2$ produce the rest of the row $(2,1,1)$.

\section{Scope}

Theorem~\ref{thm:recurrence} answers the recurrence clause in the natural
coefficient direction: for every fixed $n$, $k\mapsto f_{n,k}$ satisfies a
constructively specified scalar finite-shift recurrence. It does not assert
a low-order recurrence uniform in $n$, nor a recurrence coupling different
values of $n$. The source question asks only whether $f_{n,k}$ satisfies
``some recurrence'' and does not prescribe such stronger uniformity.
The paper's random linear-independence conjecture and Gram-matrix positivity
question are also outside the scope of this result.

\section{Exact computational verification}

The included script \texttt{code/verifier.py} implements every term in
\eqref{eq:main} with integer and rational arithmetic. Run from the repository
root:
\begin{verbatim}
conda run --no-capture-output -n sandbox python \
  runs/fa_banach_001/solutions/full/\
  2503.09542_orbit_count_cycle_index_formula/code/verifier.py
\end{verbatim}
It reproduces every source-table value through $n=5$, $k=12$:
\begin{center}
\begin{tabular}{@{}ccl@{}}
\toprule
$n$ & tested $k$ & final tested value \\
\midrule
2 & $0\leq k\leq2$ & $f_{2,2}=1$ \\
3 & $0\leq k\leq6$ & $f_{3,6}=1$ \\
4 & $0\leq k\leq12$ & $f_{4,12}=5050$ \\
5 & $0\leq k\leq12$ & $f_{5,12}=732149722382$ \\
\bottomrule
\end{tabular}
\end{center}
As a second check, the script directly enumerates the cycles of
$\sigma\mapsto\mu\sigma\nu$ for all 39 pairs of conjugacy types across
$1\leq n\leq4$ and compares them to the power-type/Moebius formula. All
comparisons agree. These tests are consistency evidence; the proof does not
depend on them.

The second script \texttt{code/recurrence\_verifier.py} computes the
dimensions of the explicit fixed-subset polynomial spans for $1\leq n\leq5$:
\[
\begin{array}{c|ccccc}
n&1&2&3&4&5\\ \hline
m_n&1&2&4&10&19
\end{array}
\]
It constructs the Wronskian operator for $n=3$ and verifies both the displayed
differential equation and every coefficient-recurrence residual exactly.
Run it with
\begin{verbatim}
conda run --no-capture-output -n sandbox python \
  runs/fa_banach_001/solutions/full/\
  2503.09542_orbit_count_cycle_index_formula/code/recurrence_verifier.py
\end{verbatim}

\section{Novelty audit and review recommendation}

The bounded search (performed 26 August 2026) covered the run registry and
solution indexes; arXiv id 2503.09542; the exact title and question wording;
searches for left-right subset orbits together with Burnside, cycle-index,
Wronskian, and recurrence
keywords; and OEIS A381842~\cite{OEIS}. The source contains no occurrence of
``Burnside'', ``Polya'', or ``cycle index''. The OEIS entry, last modified
16 July 2026, lists only $f_{n,2}=p(n)-1$, complement symmetry, and initial
values. A fresh recurrence-focused search on 26 August 2026 found no later
paper or matching formula.

Novelty confidence is nevertheless only \emph{moderate}: this is a natural
finite-group enumeration argument and could appear in older literature under
different terminology such as two-sided translations of subsets. Human review
should first verify equations~\eqref{eq:fixpower},~\eqref{eq:fixedsubset},
and the Wronskian elimination in Theorem~\ref{thm:recurrence},
then commission a specialist historical search before any public novelty
claim. The mathematical validity of both the formula and recurrence appears
high. A reviewer should also confirm that recurrence in $k$ is the intended
reading of the source's deliberately unspecified recurrence clause.

\begin{thebibliography}{9}
\bibitem{KT}
A. Kushwaha and R. Tripathi,
\emph{A note on Erd\H{o}s matrices and Marcus--Ree inequality},
Linear Algebra Appl. \textbf{725} (2025), 223--247;
\href{https://arxiv.org/abs/2503.09542}{arXiv:2503.09542}.

\bibitem{Cameron}
P. J. Cameron,
\emph{Combinatorics: Topics, Techniques, Algorithms},
Cambridge University Press, 1994.

\bibitem{Ince}
E. L. Ince,
\emph{Ordinary Differential Equations},
Dover Publications, 1956.

\bibitem{OEIS}
OEIS Foundation Inc.,
\emph{The On-Line Encyclopedia of Integer Sequences},
\href{https://oeis.org/A381842}{A381842}, accessed 26 August 2026.
\end{thebibliography}

\end{document}
